\documentclass[12pt]{article}
\usepackage[utf8]{inputenc}
\usepackage[T1]{fontenc}
\usepackage{hyperref}
\title{Gabrielino-Fernande\~no Language Resources for K-12 Curriculum Development}
\author{Compiled by Claude Sonnet 5.5 (claude-sonnet-5-5) for Andr\'as Kornai}
\date{}
\begin{document}
\maketitle

\section{Introduction}
Gabrielino-Fernande\~no, also called Tongva (ISO 639-3 xgf), is a Uto-Aztecan language of the Los Angeles basin. It is dormant, and sources describe it as ``reawakening''. The tribal council has worked on revival since about 2012 through written vocabularies and classes.

\section{Primary Educational Resources}
\begin{itemize}
\item \emph{Yaara' Shiraaw'ax '\!Eyooshiraaw'a: Now You're Speaking Our Language}, by Pamela Munro and the Gabrielino/Tongva Language Committee (2008, 2nd ed.\ 2012). It is a dictionary and lesson book compiled largely from J.\ P.\ Harrington's archival notes. A catalogue entry is at \url{https://wiki.ubc.ca/Documentation:RelLex/Yaara%27_Shriraaw%27ax_%27Eyooshiraaw%27a:_Now_You%27re_Speaking_Our_Language_Gabrielino/Tongva/Fernande%C3%B1o}
\item The Zócalo Public Square article (2014) mentions a committee phrasebook and a children's animal book: \url{https://live-zocalopublicsquare.ws.asu.edu/2014/10/27/topanga-cahuenga-and-tujunga-sounds-from-a-rediscovered-local-language/ideas/nexus}
\end{itemize}

\section{Linguistic and Reference Materials}
\begin{itemize}
\item Wikipedia, ``Tongva language'': \url{https://en.wikipedia.org/wiki/Tongva_language}
\item Omniglot, Tongva: \url{https://www.omniglot.com/writing/tongva.htm}
\end{itemize}

\section{Religious and Community Resources}
The Gabrielino-Tongva Language Committee has produced grammar lessons and songs (per the Wikipedia article). I found no committee web address in the search results. No scripture translations were found.

\section{Digital and Audio Resources}
The Language Committee is reported to post a daily audio word, phrase or song on Facebook. No URL was seen, so it is not listed in the json.

\section{Curriculum Development Context}
Tongva has a published dictionary and lessons, an active committee and some second-language learners, so it is among the better-resourced California revitalization cases. A curriculum should be built with the tribal language committee. I did not find current class schedules.
\end{document}